\chapter{Une carte de tous les marchés}\label{ch:m3:a-map-of-all-markets}

Le marché des changes échange près de dix trillions de dollars par jour~; le perpétuel bitcoin le plus actif sur
la plus grande plateforme crypto, moins de vingt milliards de dollars~; les contrats météo du \cref{ch:m3:freight-weather-and-other-corners}, bien moins. Les trois premiers
livres de cette série ont décrit des dizaines de marchés un par un, chacun avec ses propres instruments,
institutions, horaires et habitudes. Ce dernier chapitre les réunit sur une seule page~: leur taille, leur
organisation, leurs heures de négociation, qui les négocie et quels types de stratégies elles récompensent. C'est une carte destinée à une
société de trading qui décide où aller ensuite, et à un lecteur qui doit choisir lequel des livres suivants ouvrir.

\section{Taille}

\begin{definition}[Vitesse de rotation]\label{def:m3:a-map-of-all-markets:velocity}
La \emph{vitesse de rotation}\index{vitesse de rotation} d'un marché est son volume quotidien divisé par le montant de l'actif
en circulation (ou, pour les dérivés, les positions ouvertes)~: la part de l'encours qui change de mains chaque jour.
\end{definition}

La taille seule induit en erreur~: un marché peut être immense et endormi, ou petit et frénétique. Le marché des bons du Trésor américain se négocie pour environ 1.2~trillion de dollars
par jour, pour 31.8~trillions de dollars en circulation, soit une \omterm{def:m3:a-map-of-all-markets:velocity}{vitesse de rotation} de 3.8~\% par jour. Le plus grand perpétuel bitcoin d'une seule plateforme a échangé environ
17.6~milliards de dollars sur les 24 heures précédant l'après-midi du 24 septembre 2026, pour des positions ouvertes d'environ 96\,000 bitcoins, soit environ 8
milliards de dollars~: une vitesse supérieure à deux, si bien que ses positions ouvertes se renouvelaient deux fois par jour. La vitesse indique quel type de négociation un marché récompense~:
faible vitesse, la fourniture patiente de capacité bilancielle (le livre~2)~; forte vitesse, la rapidité et la rotation de l'inventaire.

\begin{dated}{2026-09}{Les tailles de marché sur une même échelle}\label{dat:m3:a-map-of-all-markets:sizes}
L'enquête triennale de la BIS évalue le volume quotidien moyen des opérations de change à 9.6~trillions de dollars en avril 2025. La SIFMA a fait état d'un volume de négociation de titres à revenu fixe aux États-Unis de 1\,659.2
milliards de dollars par jour, dont 1\,203.6~milliards de dollars pour les bons du Trésor américain, pour la période de 2026 à août, avec 31.8~trillions de dollars de bons du Trésor en circulation. La FIA a dénombré
119.29~milliards de contrats à terme et d'options négociés en bourse dans le monde en 2025. Le 24 septembre 2026, les points d'accès publics de Binance affichaient, sur 24 heures,
un volume de cotation de 17.65~milliards de dollars sur le perpétuel BTCUSDT et 1.90~milliards de dollars sur le comptant BTCUSDT, et des positions ouvertes de 96\,463 bitcoins sur le
perpétuel~; celui de Deribit affichait 1.58~milliards de dollars de volume sur 24 heures. En 2025, les contrats à terme et options du LME ont affiché en moyenne 717\,334 lots par
jour, et les marchés de l'électricité d'EPEX SPOT ont échangé 917.5~TWh sur l'année.
\end{dated}

\begin{omfigure}
\begin{tikzpicture}
\begin{axis}[width=11cm,height=5.2cm,xbar,bar width=9pt,xmode=log,log basis x=10,
  xlabel={volume quotidien, milliards de USD (échelle logarithmique)},
  tick label style={font=\small},label style={font=\small},
  ymin=0.4,ymax=6.6,ytick={1,2,3,4,5,6},
  yticklabels={change (avr. 2025),obligations américaines,bons du Trésor américains,perpétuel BTC (Binance),BTC comptant (Binance),perpétuel BTC (Deribit)},
  y dir=reverse,xmin=1,xmax=20000,grid=major,grid style={gray!25},
  table/col sep=comma]
\addplot[fill=omThm!60,draw=omThm] table[x=usd_bn,y=idx]{figdata/markets-3/29-a-map-of-all-markets/sizes.csv};
\end{axis}
\end{tikzpicture}

\omcaption{Volume quotidien de six marchés sur une même échelle logarithmique. Les mesures diffèrent (la BIS compte tous les instruments de change, nets des doubles comptages
en avril 2025~; la SIFMA moyenne les échanges américains de 2026 jusqu'à août~; les chiffres crypto correspondent au volume d'une journée sur une seule plateforme), et l'écart entre le
plus grand et le plus petit atteint près de quatre ordres de grandeur. Données~: \cref{dat:m3:a-map-of-all-markets:sizes}.}\label{fig:m3:a-map-of-all-markets:sizes}
\end{omfigure}

\section{Structure}

Chaque marché de ces livres se situe quelque part entre deux pôles. À l'un, un carnet d'ordres à cours limité central sur une bourse, ouvert à de nombreux participants, avec une
contrepartie centrale derrière chaque transaction~: actions et contrats à terme (le livre~1), enchères day-ahead de l'électricité et plateformes crypto de ce livre. À l'autre,
les dealers cotent auprès de leurs clients sur demande de cotation (RFQ), avec un crédit et un règlement-livraison bilatéraux~: la plupart des obligations, des swaps, des opérations de change à terme et des matières premières physiques (le
livre~2 et la première partie de ce livre). Entre les deux se trouvent les hybrides~: plateformes d'obligations tous-à-tous, swaps compensés négociés sur des plateformes électroniques,
marchés de matières premières intermédiés, et les \omterm{def:m3:automated-market-makers:amm}{teneurs de marché automatisés} (AMM) et \omterm{def:m3:on-chain-order-books-and-perpetual-dexs:perpdex}{carnets d'ordres on-chain} des chapitres crypto, qui suppriment à la fois le dealer et la
chambre de compensation et mettent un protocole à leur place.

La position d'un marché entre les pôles détermine qui peut le négocier et comment. Une société de trading disposant de technologie mais de peu de bilan se tourne vers les carnets d'ordres~;
une société disposant de bilan et de relations clientèle se tourne vers les marchés de dealers~; une société capable de lire un protocole se tourne vers les chaînes. La tendance des
vingt dernières années est allée vers le pôle des carnets d'ordres, électronique et compensé, mais les marchés de dealers restent parmi les plus importants en valeur, le FX en tête.

\section{Horaires}

\begin{definition}[Trading follow-the-sun]\label{def:m3:a-map-of-all-markets:sun}
Le \emph{trading follow-the-sun}\index{trading follow-the-sun} est l'organisation d'un portefeuille géré en continu par des desks situés dans des
fuseaux horaires successifs, chacun transmettant ses positions, ses ordres et ses questions en suspens au suivant à la fin de sa journée.
\end{definition}

Les marchés ont des horaires, et ces horaires comportent des moments fixes auxquels les prix sont fixés et les positions valorisées. \Cref{fig:m3:a-map-of-all-markets:day} situe quelques-uns des
moments fixes rencontrés dans ce livre un jour d'été boréal en UTC~: le financement toutes les huit heures des \omterm{def:m3:perpetual-futures:perp}{contrats perpétuels} (\cref{ch:m3:perpetual-futures}), l'échéance à 08:00
UTC des options crypto (\cref{ch:m3:dated-futures-options-and-etfs}), l'enchère day-ahead européenne à 12:00, heure d'Europe centrale
(\cref{ch:m3:power-markets-design}), la publication des \omterm{def:m3:metals:official}{prix officiels du LME} à l'issue de la seconde séance Ring (\cref{ch:m3:metals}), la séance principale des actions américaines
de 9:30 à 16:00, heure de New York, et la tranche de 15:00 à 16:00, heure de New York, à partir de laquelle est calculé le taux de référence bitcoin pour les fonds américains. Autour
de ces moments, les contrats à terme électroniques et le FX tournent presque en continu les jours ouvrés~: le marché FX de CME Group sur Globex, par exemple, se négocie
du dimanche 17:00 au vendredi 16:00, heure de Chicago, avec une pause d'une heure chaque jour à partir de 16:00.

\begin{omfigure}
\centering
\begin{tikzpicture}[font=\footnotesize,x=0.47cm,y=0.62cm]
\draw[thick,->] (0,0) -- (24.6,0) node[right] {UTC};
\foreach \h in {0,4,8,12,16,20,24} {\draw (\h,0.12) -- (\h,-0.12) node[below] {\h:00};}
\fill[omThm!25] (0,0.3) rectangle (24,0.7); \node[anchor=west] at (0.1,0.5) {crypto~: en continu, tous les jours};
\foreach \h in {0,8,16} {\draw[omThm,very thick] (\h,0.2) -- (\h,0.8);}
\node[anchor=west] at (16.1,0.95) {financement des contrats perpétuels à 00, 08, 16};
\fill[omExo!40] (8,1.3) rectangle (8.15,1.7); \node[anchor=east] at (7.9,1.5) {échéance des options crypto 08:00};
\fill[omDef!40] (10,2.0) rectangle (10.15,2.4); \node[anchor=east] at (9.9,2.2) {enchère day-ahead UE 10:00};
\fill[omProp!40] (11.33,2.7) rectangle (12.42,3.1); \node[anchor=east] at (11.2,2.9) {prix officiels du LME 11:20--12:25};
\fill[omThm!40] (13.5,3.4) rectangle (20,3.8); \node[anchor=east] at (13.4,3.6) {séance principale des actions américaines 13:30--20:00};
\fill[omExo!60] (19,4.1) rectangle (20,4.5); \node[anchor=east] at (18.9,4.3) {fenêtre du taux de référence bitcoin 19:00--20:00};
\end{tikzpicture}

\omcaption{Moments fixes d'un jour de semaine d'été boréal, en UTC, extraits des chapitres de ce livre~: négociation crypto en continu avec financement toutes les huit heures, les
échéances des options crypto, l'enchère day-ahead européenne (12:00 CET), la publication des \omterm{def:m3:metals:official}{prix officiels du LME} (12:20--13:25 Londres), la séance principale des actions américaines (9:30--16:00
heure de New York), la fenêtre du taux de référence bitcoin de New York (15:00--16:00 heure de New York). Schéma indicatif.}\label{fig:m3:a-map-of-all-markets:day}
\end{omfigure}

\begin{proposition}[Postes nécessaires pour couvrir un portefeuille négocié en continu]\label{prop:m3:a-map-of-all-markets:shifts}
Pour doter en personnel un portefeuille négocié en continu avec des postes de $L$ heures qui se chevauchent sur $o$ heures pour la passation, un desk a besoin de $\lceil 24/(L - o)\rceil$ postes par jour~;
le nombre de traders sur chaque poste est le plus grand nombre de marchés ouverts simultanément pendant ce poste.
\end{proposition}

\begin{proof}
Les postes successifs débutent à $L - o$ heures d'intervalle, donc $n$ postes couvrent $n(L - o)$ heures du cycle, et $n$ doit atteindre 24. Un trader par marché ouvert est nécessaire à la
minute la plus chargée de chaque poste.
\end{proof}

\begin{omfigure}
\begin{tikzpicture}
\begin{axis}[width=11cm,height=4.6cm,ybar=0pt,bar width=4pt,
  xlabel={heure (UTC)},ylabel={desks ouverts},
  tick label style={font=\small},label style={font=\small},
  xmin=-0.8,xmax=23.8,ymin=0,ymax=3.5,ytick={0,1,2,3},xtick={0,3,6,9,12,15,18,21},grid=major,grid style={gray!25},
  legend columns=2,legend style={font=\footnotesize,at={(0.5,-0.34)},anchor=north,draw=none,fill=none,/tikz/every even column/.append style={column sep=8pt}},
  table/col sep=comma]
\addplot[fill=omThm!50,draw=omThm] table[x=hour,y=weekday]{figdata/markets-3/29-a-map-of-all-markets/staffing.csv};\addlegendentry{jour ouvré}
\addplot[fill=omDef!40,draw=omDef] table[x=hour,y=weekend]{figdata/markets-3/29-a-map-of-all-markets/staffing.csv};\addlegendentry{week-end}
\end{axis}
\end{tikzpicture}

\omcaption{Desks qu'une société doit doter en personnel heure par heure, pour les choix du problème du week-end~: crypto en continu, électricité européenne 05:00--17:00 UTC et marchés américains
12:00--21:00 UTC en semaine. Trois équipes de neuf heures, qui se chevauchent d'une heure, couvrent le cycle (la \cref{prop:m3:a-map-of-all-markets:shifts})~; les
heures les plus chargées exigent trois traders. Horaires illustratifs des desks. Données~: le tutoriel du chapitre.}\label{fig:m3:a-map-of-all-markets:staffing}
\end{omfigure}

\section{Participants}

Les mêmes quelques catégories de sociétés reviennent sur tous les marchés de ces livres, dans des proportions variables~: les dealers (banques et leurs successeurs), qui mettent à disposition leur bilan et
prennent le contre-pied de leurs clients~; les sociétés de trading pour compte propre, qui négocient leur propre capital, surtout dans les carnets d'ordres et de plus en plus sur les marchés des dealers
(livre~1)~; les sociétés de gestion et les hedge funds, qui prennent position sur le marché et paient pour la liquidité~; les acteurs physiques, producteurs, consommateurs et \omterm{def:m3:physical-commodity-markets:house}{maisons de négoce} des
chapitres consacrés aux matières premières, dont le trading suit leurs actifs~; et, dans la crypto, un grand nombre de traders particuliers, des émetteurs de \omterm{def:m3:blockchains-for-traders:key}{jetons} et des protocoles. Qui est
de l'autre côté décide de l'edge d'une stratégie~: dans les marchés où les hedgers et les particuliers négocient pour d'autres raisons que le prix, un teneur de marché encaisse la fourchette~; dans les marchés de
professionnels uniquement, l'edge provient de l'information, de la vitesse ou du bilan.

\section{Types de stratégies dominantes}

Sur l'ensemble des marchés de ces livres, les stratégies rentables se répartissent en quelques familles, que les livres suivants reprennent une à une. Market making et fourniture de
liquidité~: actions, contrats à terme, options, change, plateformes et pools crypto (livres~10 et~11 sur la microstructure et le trading haute fréquence). Relative value et
arbitrage (livre~9)~: basis trades, \omterm{def:m3:spot-markets:cross}{arbitrage entre plateformes}, spreads entre contrats liés, allant des contrats à terme sur le Trésor à la \omterm{def:m3:spot-markets:kimchi}{prime kimchi}. Carry et primes de risque (livres~8 et~9)~: trades de
funding, \omterm{def:m3:forward-curves-storage-and-convenience-yield:roll}{rendement de roll}, vente de volatilité, \omterm{def:m3:prediction-and-betting-markets:flb}{biais favori-outsider}. Stratégies directionnelles et systématiques (livre~8)~: tendance et momentum sur les contrats à terme, value et
stratégies factorielles sur actions. Et l'optimisation physique~: stockage, transport et production d'électricité dans les matières premières, où la transaction revient à exploiter un actif. La carte d'une société est
l'intersection de ces familles avec les marchés que la société peut servir, compte tenu de leur structure, de leurs horaires et de leurs participants.

\section{Tutoriel~: la carte et le desk}\label{tut:m3:a-map-of-all-markets:map}

\begin{tutorial}
\textbf{Objectif.} Placer les tailles de marché sur une même échelle logarithmique, calculer les \omterm{def:m3:a-map-of-all-markets:velocity}{vitesses de rotation} et planifier les équipes couvrant en continu les marchés choisis par une société.
\textbf{État final~:} la \cref{fig:m3:a-map-of-all-markets:sizes} et les chiffres du problème du week-end.
\begin{tutsteps}
\item \textbf{Le planificateur.} Le plus petit nombre d'équipes se chevauchant, couvrant chaque minute dotée de personnel, et le nombre de traders par équipe.
\omcode{code/firm/marketmap/firm_marketmap.py}{40}{61}{Planification des équipes sur une journée de séance de marché.}\label{lst:m3:a-map-of-all-markets:shifts}
\item \textbf{Le desk.} \texttt{desks()} fixe les choix de couverture de la société~; \texttt{shifts(9, 1)} planifie un jour de semaine et un jour de week-end.
\item \textbf{Exécuter} \texttt{turnover\_velocity} pour les bons du Trésor et un \omterm{def:m3:perpetual-futures:perp}{contrat perpétuel}, puis \texttt{fig\_map.py}.
\end{tutsteps}
\textbf{Pour aller plus loin.} Ajouter un desk asiatique et voir s'il permet d'économiser une équipe~; prévoir les changements d'heure~; compter les heures-trader par semaine et comparer avec
le coût d'une relève de nuit automatisée, avec un opérateur d'astreinte.
\end{tutorial}

\section{Construction~: le registre des marchés}\label{bld:m3:a-map-of-all-markets:marketmap}

\begin{build}
\bfield{Rôle} La société miniature doit savoir, pour chaque marché qu'elle négocie, quand il est ouvert, comment il est structuré et compensé, et quelle est sa taille~; ses opérations
doivent planifier leurs effectifs et leurs systèmes autour de ces horaires.
\bfield{Interface} \texttt{Market(name, daily\_usd\_bn, structure, cleared, hours)}~; \texttt{every\_day(start, end, days)}~; \texttt{is\_open},
\texttt{open\_markets(markets, weekday, minute)}~; \texttt{staffed\_minutes}~; \texttt{plan\_shifts(markets, weekday, length\_h, overlap\_h)}.
\bfield{Règles} Horaires en minutes UTC par jour de semaine, convertis à partir des règles locales en tenant compte de l'heure d'été et des jours fériés (depuis \texttt{firm.calendar} dans une
implémentation complète)~; les tailles conservées avec la mesure de leur source~; un marché ouvert en permanence planifié comme un cycle.
\bfield{Tests d'acceptation} \texttt{code/firm/marketmap/tests/}~: requêtes d'ouverture en semaine et le week-end~; trois équipes de neuf heures couvrent le cycle~; une équipe couvre la
session américaine~; aucune équipe les jours où rien n'est ouvert.
\bfield{Pour aller plus loin} Changements d'heure~; calendriers de jours fériés par marché~; listes de relève générées à partir des positions ouvertes.
\end{build}

\begin{omsources}
\item BIS, Triennial Central Bank Survey 2025 (carnet du livre~2)~; SIFMA, US Treasury and US fixed income securities statistics, September 2026~;
FIA, ETD volume, December 2025 (carnet du livre~1).
\item Points d'accès publics des API de Binance et Deribit, 24 septembre 2026.
\item NYSE, horaires et calendriers~; chapitres~5, 8, 17 et~19 de ce livre pour les autres moments fixes.
\item CME Group, trading hours page (FX Spot+ on Globex, 2025~; copie Internet Archive)~; LME, data highlights 2025 (janvier 2026)~; EPEX SPOT, annual
trading results 2025 (19 janvier 2026).
\end{omsources}

\section{Exercices}

\begin{exercise}[$\star$]\label{exo:m3:a-map-of-all-markets:1}
Calculer la \omterm{def:m3:a-map-of-all-markets:velocity}{vitesse de rotation} du marché américain du Trésor à partir de l'\cref{dat:m3:a-map-of-all-markets:sizes}.
\end{exercise}

\begin{exercise}[$\star$]\label{exo:m3:a-map-of-all-markets:2}
Convertir en UTC, pendant l'été boréal, la séance cœur des actions américaines et la fenêtre du taux de référence bitcoin de New York.
\end{exercise}

\begin{exercise}[$\star$]\label{exo:m3:a-map-of-all-markets:3}
Placer trois marchés au choix entre le pôle carnet d'ordres et le pôle dealer, et justifier chaque placement.
\end{exercise}

\begin{exercise}[$\star\star$]\label{exo:m3:a-map-of-all-markets:4}
Déterminer le nombre d'équipes de huit heures, se chevauchant d'une heure, dont a besoin un desk ouvert en continu, puis celui pour des équipes de douze heures.
\end{exercise}

\begin{exercise}[$\star\star$]\label{exo:m3:a-map-of-all-markets:5}
Expliquer pourquoi la comparaison du volume traité sur le marché des changes avec celui d'une seule plateforme crypto n'a qu'une valeur indicative.
\end{exercise}

\begin{exercise}[$\star\star$]\label{exo:m3:a-map-of-all-markets:6}
Expliquer ce qu'une vitesse supérieure à deux pour un \omterm{def:m3:perpetual-futures:perp}{contrat perpétuel} dit de ceux qui le négocient.
\end{exercise}

\begin{exercise}[$\star\star\star$]\label{exo:m3:a-map-of-all-markets:7}
\emph{Programmation.} À l'aide de \texttt{shifts}, planifier les desks du tutoriel avec des équipes de neuf heures un jour de semaine et un dimanche. Déterminer le nombre de créneaux de trader requis pour chacun.
\end{exercise}

\begin{exercise}[$\star\star\star$]\label{exo:m3:a-map-of-all-markets:8}
\emph{Repérer l'erreur.} «~Les plus grands marchés sont les meilleures opportunités.~»
\end{exercise}

\section{Problème~: le desk ouvert en continu}

\begin{problem}[{Problème du week-end --- dotation en personnel pour la crypto, l'électricité européenne et les marchés américains}]\label{pb:m3:a-map-of-all-markets:1}
Une société négocie la crypto en continu, l'électricité européenne de 05:00 à 17:00 UTC et les taux et actions américains de 12:00 à 21:00 UTC en semaine. Les équipes durent neuf
heures et se chevauchent d'une heure pour la relève~; chaque marché ouvert exige un trader.

\medskip\noindent\textbf{Partie I --- Les équipes.}
\begin{enumerate}
\item Déterminer le nombre d'équipes par jour que nécessite le seul carnet crypto.
\item Préciser les heures de début et de fin des équipes si la première commence à minuit UTC.
\item Déterminer le nombre de traders que requiert chaque équipe en semaine.
\item Faire de même pour le week-end.
\item Déterminer le nombre de créneaux de trader par semaine que cela représente.
\end{enumerate}

\medskip\noindent\textbf{Partie II --- Les moments.}
\begin{enumerate}[resume]
\item Quels moments fixes de la \cref{fig:m3:a-map-of-all-markets:day} tombent dans chaque équipe~?
\item Quelle équipe prend en charge l'enchère day-ahead européenne~?
\item Quelle équipe couvre la clôture des actions américaines et le taux de référence du bitcoin~?
\item Que doit couvrir chaque relève~?
\item Comment le passage à l'heure d'hiver modifie-t-il le plan~?
\end{enumerate}

\medskip\noindent\textbf{Partie III --- Les alternatives.}
\begin{enumerate}[resume]
\item Quelles économies et quels coûts entraîneraient des équipes de douze heures~?
\item Un bureau en Asie pourrait-il remplacer l'équipe de nuit~?
\item Qu'est-ce qui peut être automatisé la nuit, et qu'est-ce qu'un humain doit surveiller~?
\item Comment assurer le week-end~?
\item Quel est le coût du plan en salaires si un trader coûte un montant annuel donné~?
\end{enumerate}

\medskip\noindent\textbf{Partie IV --- Jugement.}
\begin{enumerate}[resume]
\item Lequel des trois marchés la société devrait-elle abandonner si elle y était contrainte, et pourquoi~?
\item Que dit la carte sur le prochain axe de développement~?
\item En quoi la \omterm{def:m3:a-map-of-all-markets:velocity}{vitesse de rotation} éclaire-t-elle le choix~?
\item Énoncer le \emph{résultat nommé}~: le nombre d'équipes et le nombre de traders par équipe qui couvrent les marchés retenus avec un chevauchement d'une heure.
\item En une phrase~: à quoi sert la carte~?
\end{enumerate}
\end{problem}

\section{Questions d'entretien}

\begin{interviewq}[$\star$ \iqroles{trader}]\label{iq:m3:a-map-of-all-markets:1}
Classez le change, les titres du Trésor américain, les actions et la crypto par volume quotidien échangé, et dites ce que ce classement masque.
\end{interviewq}

\begin{interviewq}[$\star$ \iqroles{risque}]\label{iq:m3:a-map-of-all-markets:2}
Quels risques présente un book transmis d'un fuseau horaire à l'autre~?
\end{interviewq}

\begin{interviewq}[$\star\star$ \iqroles{chercheur}]\label{iq:m3:a-map-of-all-markets:3}
Comment estimeriez-vous le pool de profits accessibles d'un nouveau marché pour une société de market making~?
\end{interviewq}

\begin{interviewq}[$\star\star$ \iqroles{trader}]\label{iq:m3:a-map-of-all-markets:4}
Dans quel marché de ces livres entreriez-vous avec une petite équipe et peu de capital, et pourquoi~?
\end{interviewq}

\begin{interviewq}[$\star\star$ \iqroles{développeur}]\label{iq:m3:a-map-of-all-markets:5}
Concevez un service qui indique à tous les systèmes d'une société quels marchés sont ouverts, en tenant compte des jours fériés et de l'heure d'été.
\end{interviewq}

\begin{interviewq}[$\star\star\star$ \iqroles{chercheur, trader}]\label{iq:m3:a-map-of-all-markets:6}
Choisissez trois marchés de ces livres dont la combinaison permet à une société d'utiliser au mieux la même technologie, les mêmes personnes et le même capital.
\end{interviewq}
